\chapter{Case study: engineering Ledgerline}

\section{Context}
Ledgerline accepts small orders from a web client and an internal support tool.
The first release has one team and modest traffic. Orders are commercially
important, must be auditable, and must not be duplicated by a retry. The team
chooses a modular monolith with a relational database. This is a deliberate
starting point, not a claim that every order system should remain a monolith.

\section{Requirements and boundaries}
The initial slice contains five obligations:

\begin{enumerate}
  \item validate positive quantities, positive integer prices, and a non-empty customer ID;
  \item calculate money exactly in integer cents;
  \item persist one accepted order and its items atomically;
  \item make retried creation idempotent through a unique key and stored result;
  \item expose health and operational signals without exposing personal data.
\end{enumerate}

The system explicitly does not include payment capture, inventory reservation,
or international tax law. Excluding those features is a safety decision: their
contracts and failure modes have not yet been designed.

\section{Architecture}
\begin{center}
\begin{tikzpicture}[node distance=4mm and 4mm,>=Latex,font=\small]
  \node[draw,rounded corners,fill=codebg,minimum width=21mm,minimum height=8mm] (client) {clients};
  \node[draw,rounded corners,fill=codebg,minimum width=25mm,minimum height=8mm,right=of client] (http) {HTTP adapter};
  \node[draw,rounded corners,fill=codebg,minimum width=29mm,minimum height=8mm,right=of http] (domain) {domain + application};
  \node[draw,rounded corners,fill=codebg,minimum width=22mm,minimum height=8mm,right=of domain] (db) {relational DB};
  \node[draw,rounded corners,fill=codebg,minimum width=27mm,minimum height=8mm,below=of domain] (telemetry) {telemetry + runbook};
  \draw[->,thick,accent] (client) -- (http);
  \draw[->,thick,accent] (http) -- (domain);
  \draw[->,thick,accent] (domain) -- (db);
  \draw[->,thick,accent] (http) |- (telemetry);
  \draw[->,thick,accent] (domain) -- (telemetry);
\end{tikzpicture}
\end{center}

The reference repository separates the domain calculation from the standard
library HTTP shell. The SQL schema repeats critical domain constraints. The
container runs as an unprivileged user with a read-only filesystem and a
healthcheck. These are modest controls, but each is visible and testable.

\section{API contract}
An illustrative request is:

\begin{lstlisting}[language={},caption={The body is a data contract, not a Python object.}]
POST /orders
Idempotency-Key: 8c4a...
Content-Type: application/json

{"customer_id":"customer-1",
 "items":[{"sku":"book","quantity":2,"unit_price_cents":1250}]}
\end{lstlisting}

On success, return a stable order identifier, status, and exact money values.
For a malformed request, return a client error without writing. For a reused
key with a different payload, reject the conflict rather than silently choosing
one request. For an unknown outcome after a timeout, let the caller retry with
the key and retrieve the stored result.

\section{Reviewing the implementation}
The Python domain module is short enough to inspect line by line:

\lstinputlisting[language=Python,caption={The domain calculation has explicit validation.}]{examples/python/ledgerline.py}

The code deliberately uses integer cents and a decimal library only for the
illustrative tax calculation. In a real financial system, tax policy, rounding
jurisdiction, currency, refunds, and audit requirements need a larger domain
model. The example is not a payment ledger.

The SQL schema is an independent guard. The test runner executes it in an
in-memory SQLite database and verifies that the view agrees with the inserted
line item. A production database choice changes locking, migrations, backups,
and operational procedures; the SQL illustrates principles rather than a vendor
recommendation.

\section{Failure analysis}
\noindent\begin{tabularx}{\textwidth}{@{}p{0.22\textwidth}p{0.27\textwidth}X@{}}
\toprule
Failure & Detection & Response \\
\midrule
Client retries after lost response & Idempotency-key lookup and duplicate metric & Return stored result; alert on unusual conflict rate. \\
Database unavailable & Error ratio, dependency latency, health check & Fail closed, preserve request identity, retry with bounded backoff. \\
Invalid quantity import & Database constraint and validation error & Reject transaction; record safe diagnostic. \\
Slow dependency added later & Trace and p95/p99 latency & Deadline, isolation, capacity review, or asynchronous workflow. \\
Secret in log & Secret scanning and log review & Revoke, scrub, investigate access, add regression check. \\
\bottomrule
\end{tabularx}

Each row names a detection and a response. “Add monitoring” would be too vague;
the signal has a source and an action.

\section{What would change at larger scale}
At higher traffic, the team might add a connection pool, read replicas, queue
workers, partitioning, or independent services. Each change creates new questions:
which data is authoritative, how are duplicates reconciled, how are versions
coordinated, and how are operators alerted? The existing boundary makes these
questions easier to isolate but does not answer them automatically.

\begin{exercise}
Extend Ledgerline with a payment authorization workflow. Add states for unknown
outcomes, a reconciliation job, an outbox event, and an audit record. Explain why
your design does not promise exactly-once effects across the payment provider.
\end{exercise}

\section{Chapter summary}
A small system can demonstrate serious engineering: explicit scope, exact money,
database invariants, idempotency, operational signals, unprivileged deployment,
and failure-oriented review. Scaling the system means revisiting constraints,
not replacing reasoning with more infrastructure.
